\label{sec:conclusion}

We set out to answer a question that looks simple from the outside:
if cloud-aware autonomous scheduling is obviously better than a fixed capture plan, what reward shaping and action-space design are needed to train such a policy in simulation?
The semester work shows that the answer is not trivial.

A reproducible simulation and training platform was built to investigate that question systematically.
Fourteen experiments varied reward structure, action representation, and learning algorithm under matched baseline conditions.
Three design levers dominated the outcome: reward information at the action clock, delegating rate control to a pointing controller rather than raw torque commands, and correct implementation of the MPO trust-region update.
Once those were in place, both SAC and MPO developed selective shuttering behaviour absent in the deterministic baseline.
Reliable quantitative outperformance of that baseline on mission return was not established within the semester scope; the main contribution is a mapped design space and an experimental protocol for continuing the search.

\textbf{Outlook.}
Near-term work: extend evaluation beyond aggregate return, tighten the protocol for longer training runs, and migrate control authority toward mission-level intent while keeping a single canonical episode artifact for train, eval, and render.
Logical extensions beyond this semester include three-dimensional attitude and orbit propagation, richer terrain and real imagery ingestion, and multi-orbit scheduling.
